% Context from chapters-tex/wavelets.tex; for mathematical reference, not compilation.
	\foralls \om \in \{V,H,D\}, \quad
	W_j^\om = 
	\text{Span}\{ \psi_{j,n_1,n_2}^\om \}_{n_1,n_2}
}
where
\eq{
	\foralls \om \in \{V,H,D\}, \quad
	\psi_{j,n_1,n_2}^{\om}(x) = \frac{1}{2^j}\psi^\om\pa{ \frac{x_1-2^j n_1}{2^j}, \frac{x_2-2^j n_2}{2^j} }
}
and where the three mother wavelets are
\eq{
	\psi^H(x)=\psi(x_1)\phi(x_2),\quad
	\psi^V(x)=\phi(x_1)\psi(x_2), \qandq
	\psi^D(x)=\psi(x_1)\psi(x_2).
}
Here the span is closed in $L^2$ on an unbounded domain; on the torus, the wavelets are periodized. Figure~\ref{fig-wavelets-2d} displays examples of these wavelets.

\myfigure{
\tabquatre{
\image{wavelets}{.24}{wavelets-2d-1}&
\image{wavelets}{.24}{wavelets-2d-2}&
\image{wavelets}{.24}{wavelets-2d-3}&
\image{wavelets}{.24}{wavelets-2d-support-corrected}\\
$\psi^H$ & $\psi^V$ & $\psi^D$ & Support
}
}{ 
	2-D wavelets and a schematic support centered at $(2^j n_1,2^j n_2)$, with width $K2^j$ (right). 
}{fig-wavelets-2d}


  
\paragraph{Haar 2-D multiresolution.}

The 2-D Haar construction gives piecewise constant approximations: functions in $V_j \otimes V_j$ are constant on squares of size $2^j \times 2^j$. Figure~\ref{fig-multires-2d-haar} shows the corresponding projections of an image.

\myfigure{
\image{wavelets}{.24}{multires-2d-haar-0}
\image{wavelets}{.24}{multires-2d-haar-1}
\image{wavelets}{.24}{multires-2d-haar-2}
\image{wavelets}{.24}{multires-2d-haar-3}
}{ 
	2-D Haar approximation $P_{V_j^O}f$ for increasing $j$.  
}{fig-multires-2d-haar}

  
\paragraph{Discrete 2-D wavelet coefficients.}

As in~\eqref{eq-hyp-wav-samples}, assume that acquisition provides inner products of the continuous signal $f$ with scaling functions at scale $2^J$, with $N=2^{-J}$ samples per direction:
\eq{
	\foralls n\in \{0,\ldots,N-1\}^2, \quad a_{J,n}=\dotp{f}{(\phi_{J,n_1}\otimes\phi_{J,n_2})^P}
}
Discrete wavelet coefficients are defined as
\eq{
	\foralls \om \in \{V,H,D\}, \; \foralls J < j \leq 0, \; \foralls 0 \leq n_1,n_2 < 2^{-j}, \quad
	d_{j,n}^\om = \dotp{ f }{ \psi_{j,n}^{\om} } .
}
Here the wavelets are periodized.
 
Approximation coefficients are defined as
\eq{
	a_{j,n}=\dotp{f}{(\phi_{j,n_1}\otimes\phi_{j,n_2})^P}.
}

\myfigure{
\image{wavelets}{.3}{wavcoefs-2d-original}
\image{wavelets}{.42}{wavcoefs-2d-coefs}
}{ 
	2-D wavelet coefficients.  
}{fig-wavcoefs-2d}

Figure~\ref{fig-wavcoefs-2d} shows how the wavelet coefficients are arranged in an image of $N^2$ pixels.
 
Figure~\ref{fig-wavcoefs-example} shows other examples of wavelet decompositions.

\myfigure{
\tabtrois{
\image{wavelets}{.3}{wavcoefs-example-square}&
\image{wavelets}{.3}{wavcoefs-example-boat}&
\image{wavelets}{.3}{wavcoefs-example-mandrill}\\
\image{wavelets}{.3}{wavcoefs-example-square-coefs}&
\image{wavelets}{.3}{wavcoefs-example-boat-coefs}&
\image{wavelets}{.3}{wavcoefs-example-mandrill-coefs}
}
}{ 
	Images (top) and their wavelet coefficients (bottom). 
}{fig-wavcoefs-example}

  
\paragraph{One step of the forward 2-D transform.}

Each step separates fine-scale approximation coefficients into three detail arrays and a coarse approximation:
\eq{
	a_{j-1} \longmapsto (a_j, d_j^H, d_j^V, d_j^D).
}
As in one dimension, this map is orthogonal. It is implemented by applying the filtering and sub-sampling formula~\eqref{eq-convol-subsample} in each coordinate direction.

One first applies 1-D horizontal filtering and sub-sampling
\eqm{
	\tilde a_j &= ( a_{j-1} \star^H \bar h ) \downarrow^H 2\\
	\tilde d_j &= (a_{j-1}\star^H\bar g) \downarrow^H 2,
}
where $\star^H$ denotes convolution along the first coordinate (called horizontal here), applied to each column of the array
\eq{
	a \star^H b_{n_1,n_2} = \sum_{m_1=0}^{P-1} a_{n_1-m_1,n_2} b_{m_1}
}
Here $a\in\CC^{P\times P}$ is a matrix and $b\in\CC^P$ is a filter. The notation $\downarrow^H 2$ denotes sub-sampling in the horizontal direction
\eq{
	(a \downarrow^H 2)_{n_1,n_2} = a_{2n_1,n_2}.
}

One then applies 1-D vertical filtering and sub-sampling to $\tilde a_j$ and $\tilde d_j$ to obtain
\eqm{
	\begin{array}{ll}
	a_j &=  ( \tilde a_j \star^V \bar h ) \downarrow^V 2,\\
	d_j^V &= ( \tilde a_j \star^V \bar g ) \downarrow^V 2,
	\end{array}\qquad	
	\begin{array}{ll}
	d_j^H &= ( \tilde d_j \star^V \bar h ) \downarrow^V 2,\\
	d_j^D &= ( \tilde d_j \star^V \bar g ) \downarrow^V 2,
	\end{array}
}
where the vertical operators are defined analogously to the horizontal operators, acting on rows.

\myfigure{
\image{wavelets}{.9}{wavalgo-2d-step-fwd}
}{ 
	Forward 2-D filterbank step.  
}{fig-wavalgo-2d-step-fwd}

\myfigure{
\tabtrois{
\image{wavelets}{.24}{wavalgo-2d-0}&
\image{wavelets}{.24}{wavalgo-2d-1}&
\image{wavelets}{.24}{wavalgo-2d-2}\\
Coefficients $a_{j-1}$ & Transform on rows & Transform on columns
}
}{ 
	One step of the 2-D wavelet transform algorithm.  
}{fig-wavalgo-2d}

These two forward steps are shown in the block diagram in Figure~\ref{fig-wavalgo-2d-step-fwd}.
The updates can be performed in place, storing all coefficients in an image of $N^2$ pixels, as shown in Figure~\ref{fig-wavalgo-2d}. This produces the familiar multiscale arrangement used in Figure~\ref{fig-wavcoefs-example}.

  
\paragraph{Fast 2-D wavelet transform.}

The 2-D FWT algorithm iterates these steps through the scales:
\begin{rs}
	\item \textbf{Input:} image $f \in \CC^{N\times N}$.
	\item \textbf{Initialization:} $a_J = f$.
	\item \textbf{For} $j=J+1,\ldots,j_0$.
		\eq{
			\begin{array}{ll}
			\tilde a_j &= ( a_{j-1} \star^H \bar h ) \downarrow^H 2, \\
			\tilde d_j &= (a_{j-1}\star^H\bar g) \downarrow^H 2, \\
			a_j &=  ( \tilde a_j \star^V \bar h ) \downarrow^V 2,
			\end{array}\qquad			
			\begin{array}{ll}
			d_j^V &= ( \tilde a_j \star^V \bar g ) \downarrow^V 2,\\
			d_j^H &= ( \tilde d_j \star^V \bar h ) \downarrow^V 2,\\
			d_j^D &= ( \tilde d_j \star^V \bar g ) \downarrow^V 2.
			\end{array}
		}
	\item \textbf{Output:} the coefficients $\{d_j^\om\}_{J<j\leq j_0,\om} \cup \{ a_{j_0} \}$.
\end{rs}


  
\paragraph{Fast 2-D inverse wavelet transform.}

The inverse transform undoes the horizontal and vertical filtering steps.
The first step computes
\eqm{
	\tilde a_j &= (a_j\uparrow^V2)\star^Vh + (d_j^V\uparrow^V2)\star^Vg,\\
	\tilde d_j &= (d_j^H\uparrow^V2)\star^Vh + (d_j^D\uparrow^V2)\star^Vg,
}
where the vertical up-sampling is
\eq{
	(a \uparrow^V 2)_{n_1,n_2} = 
	\choice{
		a_{n_1,k}\qifq n_2=2k, \\
		0\qifq n_2=2k+1.
	}
}
The second inverse step computes
\eqm{
	a_{j-1} &= (\tilde a_j\uparrow^H2)\star^Hh + (\tilde d_j\uparrow^H2)\star^Hg.
}
Figure~\ref{fig-wavalgo-2d-step-bwd} gives the inverse filter-bank diagram corresponding to Figure~\ref{fig-wavalgo-2d-step-fwd}.

\myfigure{
\image{wavelets}{.9}{wavalgo-2d-step-bwd}
}{ 
	Backward 2-D filterbank step.  
}{fig-wavalgo-2d-step-bwd}

The inverse fast wavelet transform iteratively applies these el